
\documentclass[10pt]{article}

% -------------------------------------------------------------------
% packages 
% -------------------------------------------------------------------

\usepackage{geometry}
\usepackage{multicol}
\usepackage{amsmath,amsthm,amsfonts,amssymb}
\usepackage{graphicx,overpic}
\usepackage{color, colortbl}
\usepackage{hyperref}
\usepackage{listings}

% -------------------------------------------------------------------
% parameters 
% -------------------------------------------------------------------

% margins & spacing 
\geometry{margin = 0.5in}
\pagestyle{empty}
\setlength{\parindent}{0pt} 
\setlength{\parskip}{0pt}

 % box style
\lstset{ 
    basicstyle = \ttfamily,
    frame = single,
    breaklines = true,
    backgroundcolor = \color{yellow},
    columns = fullflexible,
}

% hyperlinks 
\hypersetup{
    colorlinks = true,
    linkcolor = blue,
    filecolor = magenta,      
    urlcolor = cyan,
}

% multicol parameters
\setlength{\premulticols}{2pt}
\setlength{\postmulticols}{2pt}
\setlength{\multicolsep}{2pt}
\setlength{\columnsep}{20pt}

% -------------------------------------------------------------------
% functions & environments 
% -------------------------------------------------------------------

% mono space font 
\newcommand{\code}[1]{\texttt{#1}}

% -------------------------------------------------------------------
% functions & environments 
% -------------------------------------------------------------------

\begin{document}
\begin{multicols*}{2}

\begin{center}
     \Large{Unix Shell Cheat Sheet}
\end{center}

This cheat sheet serves is a reference guide to Unix Shell and is not a substitute for a formal tutorial. If you're new to the party I highly recommend reading The Carpentries' \href{https://swcarpentry.github.io/shell-novice/}{Unix Shell Tutorial} 

\begin{center}
     \Large{Common Commands}
\end{center}

Change directory to \code{/file/path}
\begin{lstlisting}
cd /file/path 
\end{lstlisting}

List contents of working directory
\begin{lstlisting}
ls 
\end{lstlisting}

\begin{center}
    \begin{tabular}{|p{80pt}p{150pt}|}
    \hline
    \code{-l} & long list format \\
    \rowcolor{yellow}
    \code{-h} & human readable format \\
    \code{-t} & sort by time  \\
    \hline
    \end{tabular}
\end{center}
 
\end{multicols*}
\end{document}

% ---------------------------------------------------------------
% template 
% ---------------------------------------------------------------